\documentclass{article}
\usepackage[T1]{fontenc}
\usepackage{lmodern}
\usepackage{graphicx}
\usepackage{amsmath,amssymb}
\usepackage[a4paper,margin=2cm]{geometry}
\usepackage[absolute,overlay]{textpos}
\setlength{\TPHorizModule}{1mm}
\setlength{\TPVertModule}{1mm}
\usepackage[indonesian]{babel}
\usepackage{hyperref}
\setlength{\emergencystretch}{2em}
% unit: Halaman 17

\begin{document}

% === Halaman 17 (python/native) ===

17

Keamanan RSA

• Keamanan algoritma RSA terletak pada tingkat kesulitan dalam 

memfaktorkan bilangan bulat n faktor-faktor prima (p dan q), yang 

dalam hal ini n = p  q.  Semakin besar n, semakin sulit difaktorkan.

• Sekali n berhasil difaktorkan menjadi p dan q, maka


 (n) = (p – 1)(q – 1) dapat dihitung.

   Selanjutnya, karena kunci enkripsi e diumumkan (tidak rahasia), maka

   kunci dekripsi d dapat dihitung dari kekongruenen

            ed  1 (mod (n)) → d =

1+𝑘 ∅(𝑛)

𝑒

\end{document}
